\begin{thebibliography}{62}

\bibitem{abramson-hasl}
H. Abramson.
A prological definition of HASL, a purely functional language with
unification-based conditional binding expressions.
\emph{New Generation Computing}, 2(1):3--35, 1984.

\bibitem{aiello-prini}
Luigia Aiello and Gianfranco Prini.
An efficient interpreter for the lambda-calculus.
\emph{Journal of Computer and System Sciences}, 23:383--424, 1981.

\bibitem{allen-lisp}
J. Allen.
\emph{Anatomy of Lisp}.
Computer Science Series. McGraw-Hill, 1978.

\bibitem{amdahl-ibm360}
G. M. Amdahl, G. A. Blaauw, and F. P. Brooks Jr.
Architecture of the IBM system/360.
\emph{IBM Journal of Research and Development}, 8:87--101, April 1964.

\bibitem{apt-van-emden}
K. R. Apt and M. H. van Emden.
Contributions to the theory of logic programming.
\emph{Journal of the ACM}, 29(3):841--862, July 1982.

\bibitem{argo-tim}
Guy Argo.
Improving the three instruction machine.
In \emph{Functional Programming Languages and Computer Architecture}, pages
100--115, Imperial College, London, September 1989.

\bibitem{backus}
J. Backus.
Can programming be liberated from the von neumann style?
\emph{Communications of the ACM}, 21(8):613--641, August 1978.

\bibitem{barbuti-leaf}
R. Barbuti, M. Bellia, G. Levi, and M. Martelli.
LEAF: a language which integrates logic, equations, and functions.
In Doug DeGroot and Gary Lindstrom, editors, \emph{Logic Programming,
Functions, Relations, and Equations}, pages 201--238. Prentice-Hall, 1986.

\bibitem{barendregt-needed}
H. P. Barendregt, J. R. Kennaway, J. W. Klop, and M. R. Sleep.
Needed reduction and spine strategies for the lambda calculus.
Technical Report CS-R8621, Centre for Mathematics and Computer Science,
Amsterdam, May 1986.

\bibitem{berkling-epsilon}
Klaus Berkling.
Epsilon-reduction: Another view of unification.
In J. V. Woods, editor, \emph{Fifth Generation Computer Architecture}. IFIP,
1986.

\bibitem{berkling-head-order}
Klaus Berkling.
Head order reduction: A graph reduction scheme for the operational lambda
calculus.
In Fasel and Keller, editors, \emph{Proceedings of the Santa-Fe Workshop on
Graph Reduction}, pages 26--48. Springer-Verlag, September/October 1986.

\bibitem{de-bruijn}
N. G. De Bruijn.
Lambda calculus notation with nameless dummies, a tool for automatic formula
manipulation, with application to the Church-Rosser theorem.
\emph{Indagationes Mathematicae}, 34(5):381--392, 1972.

\bibitem{church}
Alonzo Church.
\emph{The Calculi of Lambda-Conversion}, volume 4 of Annals of Mathematics
Studies.
Princeton University Press, 1941.

\bibitem{clark-negation}
Keith L. Clark.
Negation as failure.
In H. Gallaire and J. Minker, editors, \emph{Logic and Data Bases}, pages
293--322. Plenum Press, New York, 1978.

\bibitem{clinger-scheme}
William Clinger et al.
The revised revised report on scheme, or an uncommon lisp.
AI Memo 848, MIT Artificial Intelligence Laboratory, August 1985.

\bibitem{cohen-prolog}
Jacques Cohen.
Describing prolog by its interpretation and compilation.
\emph{Communications of the ACM}, 28(12):1311--1324, December 1985.

\bibitem{darlington-unification}
J. Darlington, A. J. Field, and H. Pull.
The unification of functional and logic languages.
In Doug DeGroot and Gary Lindstrom, editors, \emph{Logic Programming,
Functions, Relations, and Equations}, pages 37--72. Prentice-Hall, 1986.

\bibitem{dershowitz-plaisted}
Nachum Dershowitz and David A. Plaisted.
Logic programming cum applicative programming.
In \emph{1985 Symposium on Logic Programming}, pages 54--66. IEEE, July 1985.

\bibitem{fairbairn}
J. Fairbairn.
A simple abstract machine to execute supercombinators.
In Fasel and Keller, editors, \emph{Proceedings of the Santa-Fe Workshop on
Graph Reduction}, pages 49--52. Springer-Verlag, September/October 1986.

\bibitem{fairbairn-wray}
J. Fairbairn and S. Wray.
TIM: A simple, lazy abstract machine to execute supercombinators.
In Gilles Kahn, editor, \emph{Functional Programming Languages and Computer
Architecture}, pages 34--45, Portland, OR, September 1987. Springer-Verlag.

\bibitem{friedman-wise}
D. P. Friedman and D. S. Wise.
CONS should not evaluate its arguments.
In \emph{Automata, Languages, and Programming}, pages 257--284. Edinburgh
University Press, Edinburgh, 1976.

\bibitem{gagliardi}
U. O. Gagliardi.
Report of workshop 4 -- software-related advances in computer hardware.
In \emph{Symposium on the High Cost of Software}, pages 99--120, Menlo Park,
CA, 1973. Stanford Research Institute.

\bibitem{goldberg-sharing}
B. Goldberg.
Detecting sharing of partial applications in functional languages.
In Gilles Kahn, editor, \emph{Functional Programming Languages and Computer
Architecture}, pages 408--425, Portland, OR, September 1987. Springer-Verlag.

\bibitem{henderson}
Peter Henderson.
\emph{Functional Programming: Application and Implementation}.
Prentice-Hall International Series in Computer Science. Prentice-Hall,
Englewood Cliffs, NJ, 1980.

\bibitem{hindley-seldin}
J. Roger Hindley and Jonathan P. Seldin.
\emph{Introduction to Combinators and \(\lambda\)-Calculus}, volume 1 of London
Mathematical Society Student Texts.
Cambridge University Press, 1986.

\bibitem{haskell}
Paul Hudak et al.
Report on the functional programming language Haskell: Draft proposed standard.
Research Report YALEU/DCS/RR-666, December 1988.

\bibitem{huet}
G. P. Huet.
A unification algorithm for typed \(\lambda\)-calculus.
\emph{Theoretical Computer Science}, 1:27--57, 1975.

\bibitem{hughes}
John Hughes.
Why functional programming matters.
\emph{The Computer Journal}, 32(2):98--107, 1989.

\bibitem{winter}
Damir A. Jamsek, Shiu-Kai Chin, Kevin J. Greene, and Paul R. Humenn.
WINTER: An architecture supporting a purely declarative language.
Technical Report 8810, CASE Center, Syracuse University, August 1988.

\bibitem{jones-metayer}
Simon B. Jones and Daniel Le Metayer.
Compile-time garbage collection by sharing analysis.
In \emph{Functional Programming Languages and Computer Architecture}, pages
54--74, Imperial College, London, September 1989.

\bibitem{steele-common-lisp}
Guy L. Steele Jr.
\emph{CommonLisp: The Language}.
Digital Press, 1984.

\bibitem{kluge}
W. E. Kluge.
An architecture for direct execution of reduction languages.
In \emph{International Workshop on High-Level Language Computer Architecture},
Fort Lauderdale, FL, May 1980.

\bibitem{landin}
P. J. Landin.
The mechanical evaluation of expressions.
\emph{The Computer Journal}, 6(4):308--320, January 1964.

\bibitem{levy}
J. J. Levy.
Optimal reductions in the lambda-calculus.
In Seldin and Hindley, editors, \emph{To H.B. Curry -- Essays on Combinatorial
Logic, Lambda Calculus and Formalism}. Academic Press, New York, 1980.

\bibitem{lindstrom}
Gary Lindstrom.
Functional Programming and the logical variable.
In \emph{12th ACM Symposium on Principles of Programming Languages}, pages
266--280. ACM, January 1984.

\bibitem{lloyd}
John Wylie Lloyd.
\emph{Foundations of Logic Programming}.
Symbolic Computation, Artificial Intelligence. Springer-Verlag, New York,
second edition, 1987.

\bibitem{loveland}
D. W. Loveland.
Near-horn prolog.
In J. Lassez, editor, \emph{Logic Programming: Proceedings of the Fourth
Conference}, pages 456--469. MIT Press, 1987.

\bibitem{mago}
G. A. Mago.
A cellular computer architecture for functional programming.
In \emph{IEEE Computer Society COMPCON}, pages 179--187, 1980.

\bibitem{mcclain}
Mark McClain et al.
\emph{AMD Am29300 Evaluation Board User's Guide}.
Advanced Micro Devices, Inc, June 1987.

\bibitem{nadathur-miller}
Gopalan Nadathur and Dale Miller.
An overview of \(\lambda\)prolog.
In J. Lassez, editor, \emph{Logic Programming: Proceedings of the Fifth
Conference}, pages 810--827. MIT Press, 1988.

\bibitem{naish}
L. Naish.
\emph{Negation and Control in PROLOG}, volume 238 of Lecture Notes in Computer
Science.
Springer-Verlag, 1986.

\bibitem{norman}
A. C. Norman.
Faster combinator reduction using stock hardware.
In \emph{1988 ACM Symposium on LISP and Functional Programming}, pages 235--243,
Snowbird, UT, July 1988.

\bibitem{peyton-jones}
S. L. Peyton Jones.
An investigation of the relative efficiencies of combinators and lambda
expressions.
In \emph{1982 ACM Symposium on LISP and Functional Programming}, pages 150--158.
ACM, 1982.

\bibitem{plotkin}
G. D. Plotkin.
A structural approach to operational semantics.
Technical Report DAIMI FN-19, Computer Science Department, University of
Aarhus, Denmark, 1981.

\bibitem{reddy-functional-logic}
U. S. Reddy.
Functional logic languages, part i.
In Fasel and Keller, editors, \emph{Proceedings of the Santa-Fe Workshop on
Graph Reduction}, pages 401--425. Springer-Verlag, September/October 1986.

\bibitem{reddy-narrowing}
Uday S. Reddy.
Narrowing as the operational semantics of functional languages.
In \emph{1985 Symposium on Logic Programming}, pages 138--151. IEEE, July 1985.

\bibitem{revesz}
G. Revesz.
\emph{Lambda-Calculus, Combinators, and Functional Programming}, volume 4 of
Cambridge Tracts in Theoretical Computer Science.
Cambridge University Press, 1988.

\bibitem{resolution}
J. A. Robinson.
A machine-oriented logic based on the resolution principle.
\emph{Journal of the ACM}, 12(1):23--41, January 1965.

\bibitem{super}
J. A. Robinson and Kevin J. Greene.
New generation knowledge processing: Final report on the SUPER system.
Technical Report 8707, CASE Center, Syracuse University, May 1987.

\bibitem{rosser}
J. B. Rosser.
Highlights in the history of the lambda calculus.
\emph{Annals Hist. Computing}, 6:337--349, 1984.

\bibitem{roylance}
Gerald Roylance.
Expressing mathematical subroutines constructively.
In \emph{1988 ACM Conference on Lisp and Functional Programming}, pages 8--13,
Snowbird, UT, July 1988. ACM.

\bibitem{scheevel}
Mark Scheevel.
NORMA: A graph reduction processor.
In \emph{1986 ACM Conference on Lisp and Functional Programming}, pages 212--219,
Cambridge, MA, August 1986. ACM.

\bibitem{schonfinkel}
Moses Schonfinkel.
Ueber die bausteine der mathematischen logik.
\emph{Mathematische Annalen}, 92:305--316, 1924.

\bibitem{smolka}
Gert Smolka.
FRESH: A higher-order language with unification and multiple results.
In Doug DeGroot and Gary Lindstrom, editors, \emph{Logic Programming,
Functions, Relations, and Equations}, pages 469--524. Prentice-Hall, 1986.

\bibitem{staples-robinson-workshop}
J. Staples and P. J. Robinson.
Unification of quantified terms.
In Fasel and Keller, editors, \emph{Proceedings of the Santa-Fe Workshop on
Graph Reduction}, pages 426--450. Springer-Verlag, September/October 1986.

\bibitem{qu-prolog}
J. Staples, P. J. Robinson, and R. A. Hagen.
QU-PROLOG -- an extended prolog for symbolic computation.
In \emph{Proceedings of the 11th Australian Computer Science Conference},
Brisbane, February 1988.

\bibitem{staples-robinson-jlp}
John Staples and Peter J. Robinson.
Efficient unification of quantified terms.
\emph{The Journal of Logic Programming}, 5:133--149, 1988.

\bibitem{stickel}
Mark E. Stickel.
A prolog technology theorem prover: Implementation by an extended prolog
compiler.
\emph{Journal of Automated Reasoning}, 4:353--380, 1988.

\bibitem{stoye}
William Robert Stoye.
\emph{The Implementation of Functional Languages using Custom Hardware}.
PhD thesis, Cambridge University, May 1985.

\bibitem{turner}
D. A. Turner.
A new implementation technique for applicative languages.
\emph{Software -- Practice and Experience}, 9:31--49, 1979.

\bibitem{wadsworth}
C. P. Wadsworth.
\emph{Semantics and Pragmatics of the Lambda Calculus}.
PhD thesis, Oxford University, 1971.

\bibitem{zhang-berkling}
Sining Zhang and Klaus Berkling.
The soundness and completeness of head order reduction.
Technical Report 8907, CASE Center, Syracuse University, May 1989.

\end{thebibliography}
